\section{\texorpdfstring{\(\mathfrak{sl}_2\)-modules}{sl2-modules}}

Recall the standard \(\mathfrak{sl}_2\)-triple
\[
    e=\begin{pmatrix}0&1\\0&0\end{pmatrix},\qquad
    f=\begin{pmatrix}0&0\\1&0\end{pmatrix},\qquad
    h=\begin{pmatrix}1&0\\0&-1\end{pmatrix},
\]
satisfying
\[
    [e,f]=h,\qquad [h,e]=2e,\qquad [h,f]=-2f.
\]

\begin{theorem}\label{thm:sl2-irreducible}
    Let \(F\) be an algebraically closed field of characteristic \(0\), and let \(V\) be a finite-dimensional irreducible \(\mathfrak{sl}_2(F)\)-module. Then there exist
    \[
        m\in \mathbb Z_{\ge 0},\qquad 0\neq v\in V
    \]
    such that
    \[
        hv=mv,\qquad ev=0,
    \]
    and
    \[
        \{v,fv,f^2v,\dots,f^mv\}
    \]
    is a basis of \(V\). In particular,
    \[
        \dim V=m+1,
    \]
    and the weights of \(h\) on \(V\) are
    \[
        m,m-2,m-4,\dots,-m.
    \]
    Moreover, for \(0\le i\le m\),
    \[
        h(f^iv)=(m-2i)f^iv,
    \]
    and for \(1\le i\le m\),
    \[
        e(f^iv)=i(m-i+1)f^{i-1}v.
    \]
\end{theorem}

\begin{proof}
    Choose an eigenvector \(0\neq u\in V\) of \(h\), say
    \[
        hu=\alpha u.
    \]
    Then
    \[
        h(eu)=[h,e]u+e(hu)=2eu+\alpha eu=(\alpha+2)eu,
    \]
    and similarly
    \[
        h(fu)=[h,f]u+f(hu)=-2fu+\alpha fu=(\alpha-2)fu.
    \]
    Thus \(e\) raises weights by \(2\), and \(f\) lowers weights by \(2\).

    Since \(V\) is finite-dimensional, the vectors
    \[
        u,eu,e^2u,\dots
    \]
    cannot all be nonzero. Let \(r\ge 0\) be minimal such that
    \[
        e^{r+1}u=0,
    \]
    and set
    \[
        v=e^ru.
    \]
    Then \(v\neq 0\), \(ev=0\), and \(v\) is again an eigenvector of \(h\); write
    \[
        hv=\lambda v.
    \]

    Now consider the subspace
    \[
        W=\operatorname{span}\{v,fv,f^2v,\dots\}.
    \]
    It is stable under \(f\) by definition, and it is stable under \(h\) because
    \[
        h(f^iv)=(\lambda-2i)f^iv.
    \]
    We claim that \(W\) is also stable under \(e\).

    For \(i\ge 1\),
    \[
        ef^iv=[e,f]f^{i-1}v+fef^{i-1}v
        =hf^{i-1}v+fef^{i-1}v.
    \]
    Iterating this identity, we obtain
    \[
        ef^iv=\sum_{j=0}^{i-1}f^jhf^{i-1-j}v.
    \]
    Since
    \[
        hf^{i-1-j}v=(\lambda-2(i-1-j))f^{i-1-j}v,
    \]
    it follows that
    \[
        ef^iv=\sum_{j=0}^{i-1}(\lambda-2(i-1-j))f^{i-1}v
        =i(\lambda-i+1)f^{i-1}v.
    \]
    Hence \(eW\subseteq W\), so \(W\) is a nonzero submodule of \(V\). By irreducibility of \(V\), we have
    \[
        W=V.
    \]

    Now let \(m\ge 0\) satisfy
    \[
        f^mv\neq 0,\qquad f^{m+1}v=0.
    \]
    Then
    \[
        V=\operatorname{span}\{v,fv,\dots,f^mv\}.
    \]
    These vectors are linearly independent, since they are eigenvectors of \(h\) with distinct eigenvalues
    \[
        \lambda,\lambda-2,\dots,\lambda-2m.
    \]
    Thus \(\{v,fv,\dots,f^mv\}\) is a basis of \(V\).

    Finally, applying the formula for \(ef^iv\) with \(i=m+1\), we get
    \[
        0=ef^{m+1}v=(m+1)(\lambda-m)f^mv.
    \]
    Since \(f^mv\neq 0\), we must have
    \[
        \lambda=m.
    \]
    Therefore
    \[
        hv=mv,\qquad ev=0,
    \]
    and the weights are exactly
    \[
        m,m-2,\dots,-m.
    \]
\end{proof}

\begin{corollary}\label{cor:sl2-unique}
    For each \(m\in\mathbb Z_{\ge 0}\), there exists, up to isomorphism, a unique irreducible \(\mathfrak{sl}_2(F)\)-module of dimension \(m+1\).
\end{corollary}

\begin{proof}
    By Theorem~\ref{thm:sl2-irreducible}, any such module admits a basis
    \[
        v,fv,\dots,f^mv
    \]
    with
    \[
        hv=mv,\qquad ev=0,
    \]
    and
    \[
        h(f^iv)=(m-2i)f^iv,\qquad
        e(f^iv)=i(m-i+1)f^{i-1}v,\qquad
        f(f^iv)=f^{i+1}v.
    \]
    Thus the action is completely determined by \(m\).
\end{proof}

\begin{remark}
    Once complete reducibility of finite-dimensional \(\mathfrak{sl}_2(F)\)-modules is known, one obtains the following consequences.
\end{remark}

\begin{corollary}\label{cor:h-integral}
    Let \(V\) be a finite-dimensional \(\mathfrak{sl}_2(F)\)-module. Then every eigenvalue of \(h\) on \(V\) is an integer.
\end{corollary}

\begin{proof}
    Decompose \(V\) as a direct sum of irreducible submodules. By Theorem~\ref{thm:sl2-irreducible}, each irreducible summand has highest weight \(m\in\mathbb Z_{\ge0}\), hence all of its weights are
    \[
        m,m-2,\dots,-m,
    \]
    which are integers.
\end{proof}

\begin{proposition}\label{prop:weight-q-minus-p}
    Let \(V\) be a finite-dimensional \(\mathfrak{sl}_2(F)\)-module, and let \(0\neq w\in V\) be an eigenvector of \(h\) with eigenvalue \(\alpha\). Suppose
    \[
        f^{q+1}w=0,\qquad f^qw\neq 0,
    \]
    and
    \[
        e^{p+1}w=0,\qquad e^pw\neq 0.
    \]
    Then
    \[
        \alpha=q-p.
    \]
\end{proposition}

\begin{proof}
    Write \(V\) as a direct sum of irreducible submodules:
    \[
        V=\bigoplus_j V_j.
    \]
    Then
    \[
        w=\sum_j w_j,\qquad w_j\in V_j,
    \]
    where each nonzero \(w_j\) is again an eigenvector of \(h\) of eigenvalue \(\alpha\).

    Fix a nonzero \(w_j\). By Theorem~\ref{thm:sl2-irreducible}, there exists a highest weight vector \(v_j\in V_j\) and an integer \(i_j\ge 0\) such that
    \[
        w_j=f^{i_j}v_j.
    \]
    If the highest weight of \(V_j\) is \(m_j\), then
    \[
        \alpha=m_j-2i_j.
    \]
    Set
    \[
        p_j=i_j,\qquad q_j=m_j-i_j.
    \]
    Then
    \[
        \alpha=q_j-p_j.
    \]

    Because the decomposition is direct, the integers \(p\) and \(q\) attached to \(w\) satisfy
    \[
        p=\max_j p_j,\qquad q=\max_j q_j.
    \]
    Since \(q_j-p_j=\alpha\) for every nonzero \(w_j\), we get
    \[
        q=\max_j(\alpha+p_j)=\alpha+\max_j p_j=\alpha+p,
    \]
    hence
    \[
        \alpha=q-p.
    \]
\end{proof}

\begin{proposition}\label{prop:existence-all-dimensions}
    For every \(m\in\mathbb Z_{\ge 0}\), there exists an irreducible \(\mathfrak{sl}_2(F)\)-module of dimension \(m+1\).
\end{proposition}

\begin{proof}
    Consider the polynomial algebra \(F[x,y]\). Define operators on \(F[x,y]\) by
    \[
        e=y\frac{\partial}{\partial x},\qquad
        f=x\frac{\partial}{\partial y},\qquad
        h=[e,f]=y\frac{\partial}{\partial y}-x\frac{\partial}{\partial x}.
    \]
    A direct computation shows that
    \[
        [e,f]=h,\qquad [h,e]=2e,\qquad [h,f]=-2f,
    \]
    so this defines an \(\mathfrak{sl}_2(F)\)-action on \(F[x,y]\).

    Fix \(m\ge 0\), and let
    \[
        V_m=\operatorname{span}_F\{x^iy^j\mid i+j=m\}.
    \]
    Then \(V_m\) is stable under \(e,f,h\), so \(V_m\) is an \(\mathfrak{sl}_2(F)\)-submodule. Clearly,
    \[
        \dim V_m=m+1.
    \]

    Take
    \[
        v=y^m.
    \]
    Then
    \[
        ev=0,\qquad hv=my^m=mv,
    \]
    so \(v\) is a highest weight vector of weight \(m\). Moreover,
    \[
        f^i(v)=m(m-1)\cdots(m-i+1)x^iy^{m-i}\qquad (0\le i\le m),
    \]
    so the vectors
    \[
        v,fv,\dots,f^mv
    \]
    span \(V_m\). Hence, by Corollary~\ref{cor:sl2-unique}, \(V_m\) is the irreducible module of highest weight \(m\).
\end{proof}

